\section{Introduction}
\label{sec:introduction}
Object detection on an embedded robot operates within a shared compute budget.
The same GPU may execute detection, mapping, learned navigation, and manipulation
policies. NVIDIA Jetson platforms that include a Deep Learning Accelerator
(DLA) offer another execution resource for convolutional networks. Moving a
detector's learned computation to this accelerator is therefore a useful design
option, even when the GPU provides lower isolated latency. Whether that option
improves a robot's behavior depends on the resulting accuracy, application
latency, and interference with concurrent workloads.

A deployment option must first be executable. TensorRT can request DLA
placement while allowing unsupported regions to execute on the GPU
\cite{nvidia_dla_working}. Such a build can run successfully without providing
complete offloading. Supported operator names alone are insufficient: tensor
shapes, normalization axes, precision, and physical I/O layouts also constrain
placement. A detector that is efficient on a GPU may consequently become a
sequence of DLA subgraphs and GPU kernels, with additional layout conversion
between them.

YOLO11 \cite{jocher2024yolo11} illustrates this difficulty. Its convolutional
backbone and neck coexist with spatial self-attention and a decoded prediction
interface. The attention forms matrix products that DLA rejects, while the
usual head flattens feature maps and decodes boxes through an interface that
also prevents strict compilation. Merely enabling GPU fallback leaves both
sources of incompatibility in the deployed system.

We present \emph{YOLO11-DLA}, a localized adaptation that preserves the
surrounding YOLO11 architecture and its distributional detection head. We
replace the terminal C2PSA attention block with C2DLA, which combines
input-dependent channel normalization with local convolution. We also replace
the exported head interface with three raw, static, four-dimensional prediction
maps. Box decoding, confidence filtering, and non-maximum suppression (NMS)
then execute on the host. The evaluated nano instance is denoted \model;
DLA throughout means NVIDIA's accelerator.

Our scope is complete placement of the \emph{learned graph} on one DLA core.
It does not include image preprocessing or variable-cardinality detection
post-processing, and it does not imply that CUDA services, transfers, or shared
DRAM disappear. This distinction is essential when interpreting the robotics
motivation: eliminating GPU-assigned learned layers is a verifiable property,
whereas recovering GPU capacity for another task is a system-level hypothesis.

The contribution is threefold. First, we give an implemented YOLO11 adaptation
with explicit equations, tensor interfaces, and a compiler-based acceptance
criterion. Second, we characterize strict FP16 execution against both a GPU
control and stock YOLO11n with DLA fallback, using build records, engine
inspection, timing traces, and layer profiles. Third, we connect that deployment
boundary to its host-side costs and specify the outstanding tests needed to
establish numerical correctness and utility under a robotic workload.

The current evidence establishes compatibility and its measured latency tradeoff.
The adapted graph compiles as one DLA loadable, whereas stock YOLO11n fails
strict compilation. Its median engine-plus-transfer latency is 20.7\% lower
than the stock fallback configuration, but 7.86 times its own GPU latency.
Neither energy savings nor improved concurrent-task performance are assumed
from that placement result. Missing measurements appear as numbered
placeholders so that the completed study can replace them directly.
